\section{Theoretical Background}\label{sec:theoretical-background}

\subsection*{Introduction}

This section defines the five constructs that ground the study: Enterprise
Architecture (EA), EA practitioners, EA practice, Reference Architectures (RAs),
and Generative AI (GenAI). It then introduces distributed cognition as the
theoretical lens that connects them. The theoretical background from this
section is used in the construction of the conceptual model in
the \textit{Conceptual Model} (Section~\ref{sec:conceptual-model}). The following subsections detail these
foundations:

\localtableofcontents\clearpage

\subsection{Enterprise Architecture}\label{subsec:enterprise-architecture}

Enterprise Architecture has roots in IBM's Business Systems Planning methodology
of the 1960s \parencite{kotusev_enterprise_2018} and emerged as a distinct
discipline with \textcite{zachman_zachman_1987}'s framework. After nearly four
decades, the field still lacks a shared definition.
\textcite{saint-louis_investigation_2019} found that EA means different things
to different people, with definitions varying across theoretical lenses and
organizational contexts.

\textcite{the_open_group_togaf_2022} defines EA as the practice of analyzing,
designing, planning, and implementing enterprise-level analysis in service of
business strategy. \textcite{lankhorst_enterprise_2017} presents EA as an
integrated set of principles, methods, and models for designing an enterprise's
organisational structure, business processes, information systems, and
infrastructure. The first definition treats EA as an active practice; the second
treats it as both process and product.

\textcite{saint-louis_investigation_2019} identified a useful distinction: EA
as a noun versus EA as a verb. As a noun, EA refers to deliverables, models, and
artifacts. As a verb, it refers to the practice and process of architecting.
This distinction matters for this study because GenAI may affect both what
practitioners produce and how they produce it.

The lack of definitional consensus has practical consequences.
\textcite{kotusev_enterprise_2018} notes that EA frameworks offer vague,
inconsistent, and conflicting recommendations, creating problems for
organizations trying to implement EA. Some view EA as primarily IT-focused;
others emphasize business-IT alignment or enterprise-wide transformation
\parencite{gartner_gartner_2008, ross_enterprise_2006}.

In our view, EA is best understood as the sociotechnical practice of analyzing,
designing, and governing an enterprise's structures, processes, and technology
to enable strategic alignment. This working definition bridges the noun and verb
perspectives: EA produces tangible artifacts (models, standards, principles)
through an ongoing practice of stakeholder negotiation, technical design, and
organizational adaptation
\parencite{saint-louis_investigation_2019, the_open_group_togaf_2022, lankhorst_enterprise_2017}.

\subsubsection{Historical Evolution}\label{subsubsec:ea-history}

EA's origins are commonly traced to \textcite{zachman_zachman_1987}'s ``A
Framework for Information Systems Architecture.'' This work introduced multiple
perspectives and abstraction levels for describing enterprises. Zachman
established the principle that a single complex object can be described in
multiple ways, depending on the purpose and type of description required
\parencite{zachman_zachman_1987}. Early EA was largely understood through
an engineering metaphor, with practitioners viewing themselves as technical
specialists designing optimal IT solutions \parencite{spewak_enterprise_1993}.

The discipline shifted in the 2000s. \textcite{ross_enterprise_2006}
repositioned EA as the organizing logic that aligns core business processes and
IT infrastructure with a firm's integration and standardization requirements
\parencite{ross_enterprise_2006}. This moved
EA from documentation toward strategic capability. The Open Group's TOGAF
framework, first published in 1995, gained prominence through successive
revisions \parencite{the_open_group_togaf_2022}.

More recently, \textcite{kotusev_enterprise_2018} challenged many assumptions
about EA in practice. He argues that the current conceptualization rests on
anecdotal stories and wishful thinking rather than empirical evidence. Through
case study research, he showed that successful EA practices diverge from
mainstream prescriptions. Rather than following a linear path toward maturity,
EA practices adapt to organizational contexts, with successful implementations
often bearing little resemblance to framework prescriptions.

\textcite{saint-louis_investigation_2019} similarly notes that despite decades
of development, EA still lacks common understanding. EA continues to evolve,
with practitioners and researchers still working out its boundaries and methods.

\subsubsection{EA as Sociotechnical Practice}\label{subsubsec:ea-sociotechnical}

EA is not purely technical work. \textcite{saint-louis_investigation_2019}
identified three contexts within which EA operates, each combining technical and
social dimensions:

\begin{enumerate}
      \item Technological: IT infrastructure, systems integration, and technical
            optimization
      \item Socio-Technological: stakeholder participation, communication, and
            collaborative decision-making
      \item Eco-Technological: ecosystem relationships, compliance requirements,
            and external standards
\end{enumerate}

EA practice requires integrating both technical and social considerations in
specific organizational settings. \textcite{gregor_enterprise_2007} demonstrated
through case studies that effective EA development requires stakeholder
participation, with solutions achieved through communication, negotiation, and
collaboration. \textcite{lankhorst_enterprise_2017} similarly characterizes EA
as a holistic approach to organisational design, where multiple
domains (organization, information, systems, products, processes, and
applications) intersect. This sociotechnical character is why GenAI integration
in EA is not simply a tool-adoption question.

\subsection{EA Practitioners and EA Practice}\label{subsec:ea-practitioners}

\subsubsection{EA Practitioners: Roles and Evolution}\label{subsubsec:ea-roles}

EA practitioners are the professionals who develop, maintain, and govern
enterprise architectures
\parencite{kotusev_enterprise_2018, the_open_group_togaf_2022}.
\textcite{kotusev_enterprise_2018} notes that EA artifacts are developed not
only by dedicated architects but also by business executives, business analysts,
project managers, and software developers. EA practice therefore extends beyond
formal ``Enterprise Architect'' job titles.

\textcite{saint-louis_investigation_2019} identified three practitioner
archetypes:

\begin{enumerate}
      \item \textbf{Specialists}: technical experts focused on selecting optimal
            solutions for organizational needs
            \parencite{saint-louis_investigation_2019}.
      \item \textbf{Integrators}: practitioners who prioritize stakeholder
            participation and motivation in the decision-making process
            \parencite{saint-louis_investigation_2019}. They bridge technical and
            business domains through negotiation.
      \item \textbf{Facilitators}: professionals who view IT implementation and
            stakeholder engagement as inseparable from broader organizational
            adaptation \parencite{saint-louis_investigation_2019}.
\end{enumerate}

In practice, EA practitioners shift between these roles depending on the
situation. \textcite{land_enterprise_2009} observe that effective architects
need competencies spanning technical expertise, communication, leadership, and
political acumen. \textcite{van_steenbergen_maturity_2011} shows that EA
effectiveness depends on maturity across multiple dimensions, not technical
capability alone.

Our position is that ``EA practitioners'' in this research encompasses all
professionals who contribute to the development and governance of architectural
artifacts: enterprise architects, solution architects, platform engineers, data
architects, and senior software engineers who shape RAs through their daily
work. \textcite{kotusev_enterprise_2018}'s finding of distributed authorship
supports this broader scope.

\subsubsection{EA Practice: Beyond Frameworks}\label{subsubsec:ea-practice-frameworks}

We understand EA practice, following \textcite{feldman_theorizing_2011}, as the
recurrent patterns of action through which practitioners accomplish
architectural work. These practices are sociomaterial accomplishments
\parencite{orlikowski_sociomaterial_2007}: they involve both human actors and
technological artifacts.

\textcite{kotusev_enterprise_2018}'s empirical research revealed four
disconnects between formal EA frameworks and actual practice.
\textcite{ross_enterprise_2006}, \textcite{fonstad_n_o_transforming_2006}, and
\textcite{ahlemann_strategic_2012} corroborate these findings:

\begin{enumerate}
      \item \textbf{Incremental Development}: EA artifacts evolve gradually
            through organizational learning, not all at once
            \parencite{kotusev_enterprise_2018}.
      \item \textbf{Distributed Authorship}: Business executives are actively
            involved, sometimes leaving enterprise architects in a secondary
            facilitator role
            \parencite{kotusev_enterprise_2018, ross_enterprise_2006}.
      \item \textbf{Broad Application}: EA applications extend beyond strategic
            initiatives to address local requirements, unforeseen demands, and
            unplanned deviations \parencite{fonstad_n_o_transforming_2006}.
      \item \textbf{Process Integration}: Successful EA integrates with
            strategic, project, and operations management
            \parencite{ahlemann_strategic_2012}.
\end{enumerate}

The point is that EA work happens through dynamic interactions among
practitioners, tools, and organizational structures. As
\textcite{feldman_theorizing_2011} argue, understanding practice means examining
how activities are generated and operate within different contexts and over time
\parencite[p.~1241]{feldman_theorizing_2011}.

\subsection{Reference Architectures and Their Lifecycle}\label{subsec:reference-architectures}

\textcite{cloutier_concept_2010} define a Reference Architecture (RA) as an
authoritative knowledge source for a specific domain that guides and constrains
how concrete architectures and solutions are instantiated.
\textcite{angelov_framework_2012} complement this by characterizing RAs as
generic architectures for a class of information systems, designed to provide
standardization, interoperability, and reuse. RAs serve four functions: they
capture proven architectural patterns, control complexity through abstraction,
promote common understanding among stakeholders, and reduce risk through reuse
of validated designs \parencite{cloutier_concept_2010}.

\textcite{cloutier_concept_2010} describe how RAs are built from proven
architectures of past and existing products yet provide guidance for future
implementations. RAs must therefore evolve with technological change and
organizational learning. However, the literature does not present a unified view
of the RA lifecycle. The following subsections review five perspectives, then
synthesize them into an operational definition and sensitizing lifecycle model.
This model serves as a structuring device for the Delphi study: Round~1
participants position their practices within the lifecycle phases, providing
empirical validation of its utility (Section~\ref{sec:methodology}).

\subsubsection{Perspective 1: Continuous Development Cycle}

\textcite{cloutier_concept_2010} present a development cycle showing the
relationship between existing architectures, RA patterns, and new
implementations. In their view, RAs should evolve continuously because
organizational needs, technologies, and standards change. This requires active
maintenance and refactoring \parencite{cloutier_concept_2010}. They stress that
the Chief Architect (business plus technical) owns the RA and that business
sponsorship is a prerequisite.

\subsubsection{Perspective 2: Evolution and Transformation}

\textcite{angelov_framework_2012} focus on how RAs evolve over time. Using
their classification framework (which categorizes RAs by context, goal, and
scope), they note that an RA may shift categories as its environment changes or
its organizational goal evolves. Examples of evolution paths include:

\begin{itemize}
      \item A preliminary RA transforming into a classical architecture when the
            domain matures
      \item Evolution from facilitation to standardization as organizational
            needs change
      \item Expansion from single organization scope to multiple organizations
\end{itemize}

To remain usable, an RA must be adapted to reflect such changes
\parencite{angelov_framework_2012}.

\subsubsection{Perspective 3: Creation and Application Context}

Creating RAs involves analyzing existing systems, identifying commonalities and
variations, and abstracting from implementation specifics
\parencite{angelov_framework_2012}. Application requires careful attention to
context, since the generic nature of RAs leads to less defined design and
application boundaries. \textcite{cloutier_concept_2010} describe the challenge
of making RAs accessible to multiple stakeholders (engineers, business managers,
customers) while keeping them generic enough for broad reuse yet concrete enough
to be actionable. This tension between abstraction and concreteness cannot be
resolved through technical means alone. It requires ongoing negotiation with the
stakeholders who must interpret and apply the guidance.

\subsubsection{Perspective 4: Stakeholder Engagement and Validation}

The literature emphasizes intensive stakeholder involvement throughout the RA's
existence. \textcite{ross_j_w_maturity_2004} found that developing an
architectural vision required approximately sixty meetings of the executive
team, illustrating the collaboration that architectural work demands.
\textcite{angelov_framework_2012} note that RA quality can be measured by its
success in fulfilling its goals, which requires ongoing validation through
implementation and feedback.

\subsubsection{Perspective 5: Ongoing Refinement}

Rather than following a linear lifecycle, the literature suggests RAs exist in a
state of continuous adaptation. \textcite{cloutier_concept_2010} explain that
RAs must be actively managed to ensure they are updated based on evolving
standards, emerging technology, changing stakeholder requirements, and maturing
business needs. The RA lifecycle is an ongoing process of refinement, not a
sequence of phases.

\subsubsection{Operational Definition for This Research}

The five perspectives can be summarized as follows:

\begin{enumerate}
      \item \textbf{Continuous Development Cycle}: RAs require active
            maintenance, not one-time creation
            \parencite{cloutier_concept_2010}.
      \item \textbf{Evolution and Transformation}: RAs shift across categories
            as technology, experience, and goals change
            \parencite{angelov_framework_2012}.
      \item \textbf{Creation and Application Context}: RA design must balance
            abstraction for reuse with concreteness for implementation
            \parencite{angelov_framework_2012, cloutier_concept_2010}.
      \item \textbf{Stakeholder Engagement}: Architectural work demands
            extensive multi-stakeholder collaboration
            \parencite{ross_j_w_maturity_2004, angelov_framework_2012}.
      \item \textbf{Ongoing Refinement}: RAs exist in continuous adaptation
            driven by implementation experience
            \parencite{cloutier_concept_2010}.
\end{enumerate}

RA work extends well beyond initial creation. For this study, Reference
Architecture encompasses templates, standards, patterns, governance documents,
blueprints, and domain models used as reference for downstream architectural
design, regardless of formality level. Project-specific solution architectures
fall outside this scope, unless they are reused as reference for other projects,
at which point they function as RAs.

This broad scope reflects how practitioners engage with RAs in practice. RA work
encompasses creating, reading, absorbing, contextualizing, validating,
communicating, governing, and evolving architectural artifacts.
\textcite{cloutier_concept_2010} describe a continuous development cycle in
which RAs require active organizational management.
\textcite{angelov_framework_2012} note that RAs must evolve to reflect
changes in technology, experience, and goals.

\subsubsection{Sensitizing Lifecycle Model}

The Delphi study requires participants to position their GenAI practices within
the RA lifecycle. However, the five perspectives describe lifecycle
\textit{characteristics} (continuous, iterative, stakeholder-intensive) without
converging on named phases. To provide a shared vocabulary for data collection,
we synthesize an eight-phase lifecycle model:

\begin{enumerate}
    \item Stakeholder Analysis \& Requirements Gathering
    \item Current State Analysis
    \item Pattern Identification \& Abstraction
    \item Target State Design
    \item Component Specification
    \item Documentation \& Communication
    \item Validation \& Review
    \item Evolution \& Maintenance
\end{enumerate}

Each phase maps to the perspectives discussed above.
Phases~1--2 correspond to requirements and context analysis
\parencite{angelov_framework_2012}. Phases~3--5 reflect the creation and
abstraction activities in the development cycle
\parencite{cloutier_concept_2010, angelov_framework_2012}. Phase~6 addresses
the multi-stakeholder accessibility that
\textcite{cloutier_concept_2010} identify as essential. Phase~7 captures the
quality validation that \textcite{angelov_framework_2012} emphasize. Phase~8
embodies the ongoing refinement both sources describe.

This model serves as a \textit{sensitizing device}, not a prescriptive
sequence. Practitioners may enter at any phase, iterate, skip phases, or work
across multiple phases simultaneously. Its purpose is to provide a common
vocabulary for participants to position their GenAI practices within the RA
lifecycle. The model is used as the data collection instrument in
Section~\ref{sec:methodology}, where Round~1 workbook participants map their
practices to these eight phases.

\subsection{Generative AI: Emergence and Capabilities}\label{subsec:generative-ai}

Artificial intelligence research has progressed through several waves: from
rule-based expert systems in the 1980s, through machine learning and deep
learning approaches that required curated training data, to the current
generation of large-scale models that learn from broad corpora and generate
novel outputs \parencite{kanbach_genai_2024}. Generative AI (GenAI) refers to
this latest class of AI systems, capable of generating new content (text,
images, audio, code, and video) based on patterns learned from training data
\parencite{kanbach_genai_2024}.

\textcite{vaswani_attention_2017} introduced the transformer architecture that
underlies current GenAI systems, enabling attention mechanisms that capture
global dependencies without relying on recurrence or convolutions. Large
Language Models (LLMs) are specific GenAI implementations focused on natural
language. These models are trained on vast text corpora using self-supervised
learning \parencite{brown_language_2020}. The key innovation is not information
retrieval but the ability to generate coherent, contextually appropriate text by
combining learned patterns in novel ways \parencite{bubeck_sparks_2023}.

\subsubsection{Rapid Adoption}

As documented in Section~\ref{sec:introduction}, GenAI adoption has been rapid
and broad. The public release of ChatGPT in November 2022 reached 100 million
users within two months \parencite{hu_chatgpt_2023}, and by 2025 the majority
of developers and knowledge workers use AI tools regularly. This adoption has
occurred despite limited understanding of best practices, creating what
\textcite{dellacqua_navigating_2023} describe as a ``jagged technological
frontier'' where capabilities are unevenly distributed across task types.

\subsubsection{Core Capabilities Relevant to EA Practice}

GenAI systems exhibit several capabilities relevant to architectural work:

\begin{enumerate}
      \item \textbf{Pattern Recognition and Synthesis}: GenAI can identify
            patterns in unstructured content (text, diagrams, code) and
            synthesize novel outputs that combine them
            \parencite{brynjolfsson_generative_2025}. For EA practitioners, this
            enables rapid analysis of existing architectures.
      \item \textbf{Context-Aware Generation}: Modern LLMs maintain context
            windows of hundreds of thousands of tokens, allowing extended
            interactions and iterative refinement of architectural artifacts
            \parencite{kanbach_genai_2024, banh_affordance_genai_ea_2025}.
      \item \textbf{Multi-Modal Reasoning}: Recent models process and generate
            text, code, images, and structured data
            \parencite{kanbach_genai_2024}, supporting the diverse
            representational needs of architectural work.
      \item \textbf{Knowledge Integration}: GenAI can combine information from
            multiple domains, bridging technical and business perspectives
            \parencite{banh_affordance_genai_ea_2025}.
\end{enumerate}

However, GenAI also introduces challenges: hallucination (generating plausible
but incorrect information) \parencite{zhang_sirens_2023}, lack of genuine
understanding, and potential for bias inherited from training data
\parencite{bender_dangers_2021}.

\subsubsection{GenAI Integration in Professional Practice}

Emerging research reveals patterns of GenAI integration relevant to EA practice.
\textcite{noy_experimental_2023} found that professionals using ChatGPT for
writing tasks decreased inequality between workers, with less skilled workers
benefiting more while overall productivity increased by 37\%.

\textcite{dellacqua_navigating_2023} identified two patterns of AI use among
consultants. Centaurs strategically divide tasks between themselves and AI.
Cyborgs integrate AI into their workflow for all tasks. Their research showed
that professionals develop contextual understanding about when AI helps versus
hinders.

\textcite{tankelevitch_metacognitive_2024} documented new metacognitive demands
when professionals work with AI:

\begin{enumerate}
      \item Evaluating AI-generated content for accuracy and relevance
      \item Managing the balance between efficiency and critical thinking
      \item Developing strategies for effective prompting
      \item Maintaining domain expertise despite AI assistance
\end{enumerate}

\textcite{lee_impact_2025} surveyed 319 knowledge workers using GenAI tools at
work and found that GenAI shifts critical-thinking effort from material
production to AI oversight, with workers reporting reduced effort for
information gathering but greater effort for verifying and integrating AI
outputs. Taken together, these findings show that GenAI does not merely
automate tasks. It redistributes cognitive work between human and machine,
which requires a theoretical framework to analyze.

\subsection{Distributed Cognition as Integrating Lens}\label{subsec:distributed-cognition}

\subsubsection{Distributed Cognition Theory}

Distributed cognition \parencite{hutchins_cognition_1995} provides a framework
for understanding how cognitive processes spread across the members of a social
group, the tools they use, and the environment in which activity takes place.
Rather than treating cognition as something that happens only inside individual
brains, this perspective examines how cognitive systems accomplish tasks through
the propagation of representational states across media.

\textcite{hutchins_cognition_1995} studied navigation teams and showed that the
computational work of navigation is accomplished through physical manipulation
of representations distributed across team members, instruments, charts, and
procedures. The cognitive system includes all these elements working together.
No individual component could accomplish the task alone.

\textcite{hollan_distributed_2000} extended this framework to human-computer
interaction, identifying three types of cognitive distribution:

\begin{enumerate}
      \item Cognitive processes distributed across group members
      \item Cognitive processes involving coordination between internal and
            external structures
      \item Processes distributed through time, with products of earlier events
            transforming later events
\end{enumerate}

\subsubsection{Application to EA-GenAI Integration}

\textcite{hutchins_cognition_1995} canonical study of ship navigation analyzed
the cognitive system as distributed across four substrate categories: team
members, instruments, external representations, and the procedures coordinating
their use. We adopt this same decomposition for the EA-GenAI system. Each
element below maps to a substrate category recognised in the distributed
cognition literature \parencite{hollan_distributed_2000,
perry_m_distributed_2003}.

In GenAI-augmented architectural work, cognitive resources are distributed
across four elements:

\begin{enumerate}
      \item \textbf{EA Practitioners}: domain knowledge, organizational
            understanding, stakeholder awareness, and evaluative judgment
            \parencite{perry_m_distributed_2003}
      \item \textbf{GenAI Systems}: pattern recognition, content generation,
            rapid synthesis, and access to encoded knowledge
            \parencite{brynjolfsson_generative_2025}
      \item \textbf{Organizational Artifacts}: external memory and constraint
            structures \parencite{zhang_representations_1994}
      \item \textbf{Collaborative Tools}: coordination among human and
            artificial agents \parencite{brynjolfsson_generative_2025}
\end{enumerate}

A clarification: the three types of distribution from
\textcite{hollan_distributed_2000} describe \emph{how} cognition distributes
(across people, between internal and external structures, and through time).
The four elements above identify \emph{across what} cognitive resources are
distributed in the EA-GenAI system. This decomposition serves as the analytical
lens through which Sections~\ref{sec:results} and~\ref{sec:discussion} examine
how introducing GenAI restructures the cognitive system for RA development. The
Delphi study empirically tests whether this theoretical mapping holds.

As \textcite{hutchins_cognition_1995} argues, the functional properties of a
cognitive system cannot be reduced to the properties of individual agents. The
distributed system may produce capabilities neither human nor AI possesses
alone. In our view, this explains why simple tool-adoption perspectives miss
what is happening when GenAI enters EA work. The unit of analysis must be the
larger system \parencite{hollan_distributed_2000}, not the individual
practitioner.

\subsection{Core Constructs for This Study}\label{subsec:coreelements}

This study examines the interaction of five constructs:

\begin{enumerate}
      \item \textbf{Generative AI (GenAI)}: the technological capabilities
            introduced into the EA ecosystem, serving as a cognitive partner
            within a distributed system
            \parencite{hutchins_cognition_1995, vaswani_attention_2017}.
      \item \textbf{EA Practitioners}: professionals who operate as specialists,
            integrators, or facilitators
            \parencite{saint-louis_investigation_2019}, spanning roles from
            enterprise architect to platform engineer.
      \item \textbf{Reference Architecture}: the focal artifact for studying
            practice evolution, capturing accumulated architectural knowledge
            \parencite{cloutier_concept_2010}.
      \item \textbf{EA Practice}: the recurrent patterns of action
            \parencite{feldman_theorizing_2011} within which architectural work
            is accomplished.
      \item \textbf{Guidance (Research Output)}: the prescriptive output that
            synthesizes findings into actionable recommendations, validated
            through panel deliberation.
\end{enumerate}

These five constructs and their interrelationships form the basis of the
conceptual model in the following section.
